% abtex2-modelo-artigo.tex, v-1.9.2 laurocesar
% Copyright 2012-2014 by abnTeX2 group at http://abntex2.googlecode.com/ 
%

% ------------------------------------------------------------------------
% ------------------------------------------------------------------------
% abnTeX2: Modelo de Artigo Acadêmico em conformidade com
% ABNT NBR 6022:2003: Informação e documentação - Artigo em publicação 
% periódica científica impressa - Apresentação
% ------------------------------------------------------------------------
% ------------------------------------------------------------------------

\documentclass[
	% -- opções da classe memoir --
	article,			% indica que é um artigo acadêmico
	11pt,				% tamanho da fonte
	oneside,			% para impressão apenas no verso. Oposto a twoside
	a4paper,			% tamanho do papel. 
	% -- opções da classe abntex2 --
	%chapter=TITLE,		% títulos de capítulos convertidos em letras maiúsculas
	%section=TITLE,		% títulos de seções convertidos em letras maiúsculas
	%subsection=TITLE,	% títulos de subseções convertidos em letras maiúsculas
	%subsubsection=TITLE % títulos de subsubseções convertidos em letras maiúsculas
	% -- opções do pacote babel --
	english,			% idioma adicional para hifenização
	brazil,				% o último idioma é o principal do documento
	sumario=tradicional
	]{abntex2}


% ---
% PACOTES
% ---

% ---
% Pacotes fundamentais 
% ---
\usepackage{lmodern}			% Usa a fonte Latin Modern
\usepackage[T1]{fontenc}		% Selecao de codigos de fonte.
\usepackage[utf8]{inputenc}		% Codificacao do documento (conversão automática dos acentos)
\usepackage{indentfirst}		% Indenta o primeiro parágrafo de cada seção.
\usepackage{nomencl} 			% Lista de simbolos
\usepackage{color}				% Controle das cores
\usepackage{graphicx}			% Inclusão de gráficos
\usepackage{microtype} 			% para melhorias de justificação
% ---
		
% ---
% Pacotes adicionais, usados apenas no âmbito do Modelo Canônico do abnteX2
% ---
\usepackage{lipsum}				% para geração de dummy text
% ---
		
% ---
% Pacotes de citações
% ---
\usepackage[brazilian,hyperpageref]{backref}	 % Paginas com as citações na bibl
\usepackage[alf]{abntex2cite}	% Citações padrão ABNT
% ---

\usepackage{tcolorbox}
\usepackage{url}

% ---
% Configurações do pacote backref
% Usado sem a opção hyperpageref de backref
\renewcommand{\backrefpagesname}{Citado na(s) página(s):~}
% Texto padrão antes do número das páginas
\renewcommand{\backref}{}
% Define os textos da citação
\renewcommand*{\backrefalt}[4]{
	\ifcase #1 %
		Nenhuma citação no texto.%
	\or
		Citado na página #2.%
	\else
		Citado #1 vezes nas páginas #2.%
	\fi}%
% ---

% ---
% Informações de dados para CAPA e FOLHA DE ROSTO
% ---
% ---
\titulo{Avaliação experimental de uma arquitetura IoT de baixo custo para telemetria veicular e segurança ativa baseada em Geocercas}
\tituloestrangeiro{Experimental evaluation of a low-cost IoT architecture for vehicle telemetry and active security based on geofences.}
\autor{WILLIAMS DAVID ESTUMANO DUARTE \\ \protect\normalfont\small Universidade Federal do Pará -- Campus Universitário de Tucuruí (CAMTUC) \\ Núcleo de Desenvolvimento Amazônico em Engenharia (NDAE)}
\local{Tucuruí, PA}
\data{2026}


% ---

% ---
% Configurações de aparência do PDF final

% alterando o aspecto da cor azul
\definecolor{blue}{RGB}{41,5,195}

% informações do PDF
\makeatletter
\hypersetup{
     	%pagebackref=true,
		pdftitle={\@title}, 
		pdfauthor={\@author},
    	pdfsubject={Modelo de artigo científico com abnTeX2},
	    pdfcreator={LaTeX with abnTeX2},
		pdfkeywords={abnt}{latex}{abntex}{abntex2}{atigo científico}, 
		colorlinks=true,       		% false: boxed links; true: colored links
    	linkcolor=blue,          	% color of internal links
    	citecolor=blue,        		% color of links to bibliography
    	filecolor=magenta,      		% color of file links
		urlcolor=blue,
		bookmarksdepth=4
}
\makeatother
% --- 

% ---
% compila o indice
% ---
\makeindex
% ---

% ---
% Altera as margens padrões
% ---
\setlrmarginsandblock{3cm}{3cm}{*}
\setulmarginsandblock{3cm}{3cm}{*}
\checkandfixthelayout
% ---

% --- 
% Espaçamentos entre linhas e parágrafos 
% --- 

% O tamanho do parágrafo é dado por:
\setlength{\parindent}{1.3cm}

% Controle do espaçamento entre um parágrafo e outro:
\setlength{\parskip}{0.2cm}  % tente também \onelineskip

% Espaçamento simples
\SingleSpacing

% ----
% Início do documento
% ----
\begin{document}

% Retira espaço extra obsoleto entre as frases.
\frenchspacing

% ----------------------------------------------------------
% ELEMENTOS PRÉ-TEXTUAIS
% ----------------------------------------------------------

%---
%
% Se desejar escrever o artigo em duas colunas, descomente a linha abaixo
% e a linha com o texto ``FIM DE ARTIGO EM DUAS COLUNAS''.
% \twocolumn[    		% INICIO DE ARTIGO EM DUAS COLUNAS
%
%---
% página de titulo
\maketitle

% resumo em português
\begin{resumoumacoluna}
   O avanço da Internet das Coisas (IoT) tem transformado significativamente o cenário da telemetria veicular e da gestão de frotas, permitindo uma transição de sistemas de rastreamento passivos para soluções de monitoramento em tempo real altamente integradas. No entanto, o desenvolvimento de sistemas que combinam segurança ativa (como geocercas e interrupção remota de ignição) com viabilidade econômica ainda enfrenta desafios técnicos e de padronização. Muitas soluções comerciais atuais possuem custos elevados, o que motiva a pesquisa por arquiteturas de hardware de baixo custo baseadas em microcontroladores acessíveis, como o ESP32, Arduino e Raspberry Pi. Além da questão do hardware, a eficácia desses sistemas depende de estratégias de comunicação robustas, que variam desde o uso de redes móveis 4G/LTE até protocolos de longo alcance e baixo consumo, como o LoRaWAN. Este trabalho apresenta uma revisão da literatura com o objetivo de analisar o estado da arte desses sistemas, identificando as arquiteturas predominantes, as estratégias de atuação e como os dados são integrados em plataformas externas. A análise busca mapear as limitações técnicas e as lacunas de segurança, fornecendo uma base sólida para o desenvolvimento de novas soluções de telemetria que sejam, ao mesmo tempo, acessíveis e eficientes na proteção de ativos veiculares.

   \vspace{\onelineskip}

   \noindent
   \textbf{Palavras-chaves}: Telemetria Veicular. IoT. Gestão de frotas. Segurança. Geocerca. Microcontrolador.
\end{resumoumacoluna}

% ---
% RESUMO EM INGLÊS (Abstract)
% ---
\vspace{2cm}
\renewcommand{\resumoname}{Abstract}
\begin{resumoumacoluna}
   \begin{otherlanguage*}{english}
      The advancement of the Internet of Things (IoT) has significantly transformed the landscape of vehicle telemetry and fleet management, enabling a transition from passive tracking systems to highly integrated real-time monitoring solutions. However, developing systems that combine active security (such as geofences and remote ignition cutoff) with economic viability still faces technical and standardization challenges. Many current commercial solutions have high costs, which motivates research into low-cost hardware architectures based on accessible microcontrollers such as the ESP32, Arduino, and Raspberry Pi. Beyond hardware issues, the effectiveness of these systems depends on robust communication strategies, ranging from the use of 4G/LTE mobile networks to long-range, low-power protocols like LoRaWAN. This paper presents a literature review aimed at analyzing the state of the art of these systems, identifying predominant architectures, actuation strategies, and how data is integrated into external platforms. The analysis seeks to map technical limitations and security gaps, providing a solid foundation for developing new telemetry solutions that are both affordable and efficient in protecting vehicle assets.

      \vspace{\onelineskip}

      \noindent
      \textbf{Keywords}: Vehicle Telemetry. IoT. Fleet Management. Security. Geofencing. Microcontroller.
   \end{otherlanguage*}
\end{resumoumacoluna}

% ]  				% FIM DE ARTIGO EM DUAS COLUNAS
% ---

% ----------------------------------------------------------
% ELEMENTOS TEXTUAIS
% ----------------------------------------------------------
\textual

% ----------------------------------------------------------
% Introdução
% ----------------------------------------------------------
\section*{Introdução}
\addcontentsline{toc}{section}{Introdução}

A crescente integração de tecnologias digitais aos sistemas de transporte tem impulsionado o desenvolvimento de soluções baseadas em Internet das Coisas (IoT), especialmente no contexto da telemetria veicular e da gestão de frotas. Nesse cenário, dispositivos embarcados capazes de coletar, processar e transmitir dados em tempo real têm se tornado fundamentais para melhorar a eficiência logística, a segurança operacional e o monitoramento de ativos.

Apesar desses avanços, ainda existem desafios relevantes na implementação dessas soluções, especialmente em relação ao custo, à padronização das arquiteturas e à confiabilidade dos mecanismos de segurança. Muitas plataformas comerciais oferecem funcionalidades avançadas, porém com alto custo de aquisição e manutenção, o que limita sua adoção em contextos de menor escala ou em países em desenvolvimento. Como alternativa, tem-se observado um crescimento no uso de arquiteturas baseadas em microcontroladores de baixo custo, como ESP32, Arduino e Raspberry Pi, que viabilizam o desenvolvimento de sistemas mais acessíveis.

Além da escolha do hardware, a efetividade desses sistemas está diretamente relacionada às estratégias de comunicação e integração de dados, envolvendo desde redes celulares (como 4G/LTE) até tecnologias de longo alcance e baixo consumo, como LoRaWAN. Paralelamente, funcionalidades de segurança ativa, como geocercas, alertas automatizados e controle remoto de ignição, têm ganhado destaque, mas ainda carecem de maior consolidação em termos de implementação e validação prática.

Diante desse contexto, este trabalho apresenta uma revisão da literatura com o objetivo de analisar o estado da arte de sistemas de telemetria veicular baseados em IoT, com foco em soluções de baixo custo e mecanismos de segurança ativa. A revisão busca identificar as principais arquiteturas de hardware adotadas, as estratégias de monitoramento e atuação implementadas, bem como as formas de integração dos dados em plataformas externas. A partir dessa análise, pretende-se evidenciar limitações técnicas recorrentes e lacunas existentes na literatura, contribuindo para o desenvolvimento de soluções mais eficientes e acessíveis.

\section{Procedimentos Metodológicos}

Esta pesquisa caracteriza-se como uma revisão de literatura de natureza qualitativa, com o objetivo de analisar o estado da arte de sistemas de telemetria veicular baseados em Internet das Coisas (IoT), com ênfase em arquiteturas de baixo custo e mecanismos de segurança ativa. Para garantir a organização e a reprodutibilidade do estudo, foi elaborado um protocolo de revisão contendo a definição do tema, os objetivos, as fontes de busca, os critérios de seleção dos trabalhos e as estratégias de extração e análise dos dados.

\subsection{Formulação das Perguntas de Pesquisa}

Com o intuito de direcionar a análise e garantir a consistência da revisão, foram definidas as seguintes perguntas de pesquisa:

\begin{itemize}
   \item Quais são as arquiteturas de hardware predominantes em sistemas de telemetria veicular de baixo custo?
   \item Quais estratégias de atuação são utilizadas para implementar mecanismos de segurança ativa, como controle de ignição e geração de alertas automáticos?
   \item Como os dados de telemetria são integrados, processados e visualizados em plataformas externas, como sistemas em nuvem, aplicações web e interfaces móveis?
\end{itemize}

\subsection{Estratégia de Busca}

A estratégia de busca foi definida com o objetivo de identificar estudos relevantes relacionados à telemetria veicular baseada em Internet das Coisas (IoT), com foco em arquiteturas de baixo custo e mecanismos de segurança. As buscas foram realizadas nas bases de dados Google Acadêmico e IEEE Xplore, considerando publicações no período de 2020 a 2026.

Inicialmente, foi elaborada uma string de busca abrangente, contendo palavras-chave e sinônimos em inglês relacionados aos principais conceitos do estudo. No entanto, verificou-se que o uso de uma string muito extensa resultava em limitações nas plataformas, como ausência de resultados relevantes ou falhas na execução da busca.

Dessa forma, as strings foram adaptadas para cada base de dados, buscando equilibrar abrangência e precisão dos resultados.

As strings completas são apresentadas a seguir para fins de reprodutibilidade do estudo.

\begin{tcolorbox}[title=String de busca — Google Acadêmico]
   ("vehicle telemetry" OR "vehicle monitoring" OR "vehicle tracking" OR
   "fleet management systems")
   AND ("fleet management" OR "vehicle management" OR "fleet monitoring"
   OR "fleet operations")
   AND ("internet of things" OR "iot" OR "machine-to-machine" OR "m2m")
   AND ("vehicle" OR "vehicular" OR "automotive" OR "fleet")
   AND ("security")
   AND ("tracking" OR "localization" OR "geolocation" OR "positioning")
   AND ("speed limiting" OR "speed control" OR "speed restriction"
   OR "speed regulation" OR "speed violation")
   AND ("microcontroller" OR "embedded system" OR "microprocessor"
   OR "hardware")
\end{tcolorbox}

\begin{tcolorbox}[title=String de busca — IEEE Xplore]
   ("vehicle telemetry" OR "vehicle monitoring" OR "vehicle tracking" OR
   "fleet management systems")
   AND ("fleet management" OR "vehicle management" OR "fleet monitoring"
   OR "fleet operations")
   AND ("internet of things" OR "iot" OR "machine-to-machine" OR "m2m")
   AND ("vehicle" OR "vehicular" OR "automotive" OR "fleet")
   AND ("security")
\end{tcolorbox}

Essa adaptação foi necessária para contornar limitações técnicas das plataformas e garantir a obtenção de resultados mais alinhados ao escopo da pesquisa.

\subsection{Seleção dos Estudos}

O processo de seleção dos estudos foi realizado em duas etapas. Inicialmente, foi feita uma triagem preliminar dos trabalhos encontrados nas bases de dados, com base na leitura dos títulos e resumos, com o objetivo de identificar estudos potencialmente relevantes ao tema da pesquisa.

Em seguida, os trabalhos pré-selecionados foram analisados de forma mais detalhada, sendo submetidos a critérios de inclusão e exclusão previamente definidos. Essa etapa teve como objetivo garantir que apenas estudos alinhados ao escopo da revisão, especialmente aqueles que envolvem o desenvolvimento de sistemas de telemetria com foco em hardware, comunicação de dados e segurança, fossem considerados na análise final.

Os critérios utilizados estão apresentados na Tabela \ref{tab:criterios_selecao}.

\begin{table}[h]
   \centering
   \caption{Critérios de inclusão e exclusão dos estudos}
   \label{tab:criterios_selecao}
   \begin{tabular}{|p{10cm}|c|}
      \hline
      \textbf{Critério}                                                                                          & \textbf{Tipo} \\
      \hline
      O trabalho propõe ou desenvolve um sistema de telemetria.                                                  & Inclusão      \\
      \hline
      O trabalho inclui a segurança (ativa, passiva ou do veículo) como um de seus objetivos ou funcionalidades. & Inclusão      \\
      \hline
      O trabalho descreve o desenvolvimento de um dispositivo físico (hardware) de telemetria.                   & Inclusão      \\
      \hline
      O trabalho não especifica minimamente os componentes técnicos (controladores, módulos ou sensores).        & Exclusão      \\
      \hline
      O sistema/dispositivo não possui a função de rastreamento (geolocalização).                                & Exclusão      \\
      \hline
      O sistema não realiza a transmissão de dados para visualização ou processamento.                           & Exclusão      \\
      \hline
   \end{tabular}
\end{table}

Ao final desse processo de seleção, foram incluídos 18 artigos, considerados aderentes ao escopo da pesquisa e, portanto, selecionados para a etapa de análise.

\subsection{Avaliação dos Estudos}

Após a seleção dos trabalhos, foi realizada a etapa de avaliação dos estudos, com o objetivo de analisar a qualidade e a profundidade técnica das propostas apresentadas. Para isso, foram definidos critérios de avaliação baseados em aspectos relevantes ao contexto de sistemas de telemetria veicular, especialmente no que diz respeito à implementação prática, comunicação de dados e interação com o veículo.

Os critérios de qualidade adotados estão apresentados na Tabela \ref{tab:criterios_qualidade}.

\begin{table}[h]
   \centering
   \caption{Critérios de avaliação da qualidade dos estudos}
   \label{tab:criterios_qualidade}
   \begin{tabular}{|p{10cm}|c|}
      \hline
      \textbf{Critério}                                                                               & \textbf{Pontuação} \\
      \hline
      Nível de detalhamento da metodologia do desenvolvimento do hardware.                            & 0 a 2              \\
      \hline
      Nível de complexidade da infraestrutura de transmissão externa de dados.                        & 0 a 2              \\
      \hline
      Nível de detalhamento dos testes realizados.                                                    & 0 a 2              \\
      \hline
      Complexidade do sistema como um todo.                                                           & 0 a 2              \\
      \hline
      O sistema interage ativamente no comportamento do motorista ou veículo, independente da função. & 0 a 2              \\
      \hline
   \end{tabular}
\end{table}

A pontuação dos critérios foi atribuída de forma inteira, sem o uso de valores decimais, variando de 0 a 2 para cada item, conforme o nível de atendimento observado em cada estudo.

A partir da aplicação desses critérios, foi possível estabelecer uma análise comparativa entre os trabalhos selecionados. Observou-se que a nota mínima obtida entre os artigos foi 7, enquanto a maioria dos estudos apresentou pontuação 9, indicando um bom nível geral de qualidade. Destaca-se ainda que dois artigos atingiram a pontuação máxima de 10, evidenciando um alto grau de detalhamento técnico, robustez experimental e complexidade de implementação.


\subsection{Considerações Finais dos Procedimentos Metodológicos}

A condução desta revisão de literatura concentrou-se nas bases de dados Google Acadêmico e IEEE Xplore. A escolha por essas plataformas deve-se, principalmente, à limitação de trabalhos diretamente relacionados ao tema específico desta pesquisa, especialmente no que se refere à integração entre telemetria veicular, baixo custo e mecanismos de segurança ativa.

Diante desse cenário, os critérios de inclusão e exclusão, bem como os critérios de avaliação da qualidade dos estudos, foram aplicados de forma a garantir a seleção de trabalhos relevantes, ainda que com um nível de rigor moderado, visando preservar um conjunto de dados suficientemente representativo para análise. Essa abordagem permitiu equilibrar a abrangência da revisão com a disponibilidade limitada de estudos diretamente aderentes ao escopo proposto.

Ressalta-se que todo o processo metodológico foi previamente estruturado por meio de um protocolo de revisão de literatura, no qual estão detalhadas as estratégias de busca, incluindo palavras-chave utilizadas, combinações de termos, lista de trabalhos encontrados, links de acesso e demais informações complementares.\footnote{Protocolo de revisão de literatura disponível em: \url{https://drive.google.com/file/d/1ycxy1grOSLdAdATd5NtJrtUyO6gxBmgF/view?usp=sharing}}

% ----------------------------------------------------------
% Seção de explicações
% ----------------------------------------------------------
\section{Exemplos de comandos}

\subsection{Margens}

A norma ABNT NBR 6022:2003 não estabelece uma margem específica a ser utilizada
no artigo científico. Dessa maneira, caso deseje alterar as margens, utilize os
comandos abaixo:

\begin{verbatim}
   \setlrmarginsandblock{3cm}{3cm}{*}
   \setulmarginsandblock{3cm}{3cm}{*}
   \checkandfixthelayout
\end{verbatim}

\subsection{Duas colunas}

É comum que artigos científicos sejam escritos em duas colunas. Para isso,
adicione a opção \texttt{twocolumn} à classe do documento, como no exemplo:

\begin{verbatim}
   \documentclass[article,11pt,oneside,a4paper,twocolumn]{abntex2}
\end{verbatim}

É possível indicar pontos do texto que se deseja manter em apenas uma coluna,
geralmente o título e os resumos. Os resumos em única coluna em documentos com
a opção \texttt{twocolumn} devem ser escritos no ambiente
\texttt{resumoumacoluna}:

\begin{verbatim}
   \twocolumn[              % INICIO DE ARTIGO EM DUAS COLUNAS

     \maketitle             % pagina de titulo

     \renewcommand{\resumoname}{Nome do resumo}
     \begin{resumoumacoluna}
        Texto do resumo.
      
        \vspace{\onelineskip}
 
        \noindent
        \textbf{Palavras-chaves}: latex. abntex. editoração de texto.
     \end{resumoumacoluna}
   
   ]                        % FIM DE ARTIGO EM DUAS COLUNAS
\end{verbatim}

\subsection{Recuo do ambiente \texttt{citacao}}

Na produção de artigos (opção \texttt{article}), pode ser útil alterar o recuo
do ambiente \texttt{citacao}. Nesse caso, utilize o comando:

\begin{verbatim}
   \setlength{\ABNTEXcitacaorecuo}{1.8cm}
\end{verbatim}

Quando um documento é produzido com a opção \texttt{twocolumn}, a classe
\textsf{abntex2} automaticamente altera o recuo padrão de 4 cm, definido pela
ABNT NBR 10520:2002 seção 5.3, para 1.8 cm.

\section{Cabeçalhos e rodapés customizados}

Diferentes estilos de cabeçalhos e rodapés podem ser criados usando os
recursos padrões do \textsf{memoir}.

Um estilo próprio de cabeçalhos e rodapés pode ser diferente para páginas pares
e ímpares. Observe que a diferenciação entre páginas pares e ímpares só é
utilizada se a opção \texttt{twoside} da classe \textsf{abntex2} for utilizado.
Caso contrário, apenas o cabeçalho padrão da página par (\emph{even}) é usado.

Veja o exemplo abaixo cria um estilo chamado \texttt{meuestilo}. O código deve
ser inserido no preâmbulo do documento.

\begin{verbatim}
%%criar um novo estilo de cabeçalhos e rodapés
\makepagestyle{meuestilo}
  %%cabeçalhos
  \makeevenhead{meuestilo} %%pagina par
     {topo par à esquerda}
     {centro \thepage}
     {direita}
  \makeoddhead{meuestilo} %%pagina ímpar ou com oneside
     {topo ímpar/oneside à esquerda}
     {centro\thepage}
     {direita}
  \makeheadrule{meuestilo}{\textwidth}{\normalrulethickness} %linha
  %% rodapé
  \makeevenfoot{meuestilo}
     {rodapé par à esquerda} %%pagina par
     {centro \thepage}
     {direita} 
  \makeoddfoot{meuestilo} %%pagina ímpar ou com oneside
     {rodapé ímpar/onside à esquerda}
     {centro \thepage}
     {direita}
\end{verbatim}

Para usar o estilo criado, use o comando abaixo imediatamente após um dos
comandos de divisão do documento. Por exemplo:

\begin{verbatim}
   \begin{document}
     %%usar o estilo criado na primeira página do artigo:
     \pretextual
     \pagestyle{meuestilo}
     
     \maketitle
     ...
     
     %%usar o estilo criado nas páginas textuais
     \textual
     \pagestyle{meuestilo}
     
     \chapter{Novo capítulo}
     ...
   \end{document}  
\end{verbatim}

Outras informações sobre cabeçalhos e rodapés estão disponíveis na seção 7.3 do
manual do \textsf{memoir} \cite{memoir}.

\section{Mais exemplos no Modelo Canônico de Trabalhos Acadêmicos}

Este modelo de artigo é limitado em número de exemplos de comandos, pois são
apresentados exclusivamente comandos diretamente relacionados com a produção de
artigos.

Para exemplos adicionais de \abnTeX\ e \LaTeX, como inclusão de figuras,
fórmulas matemáticas, citações, e outros, consulte o documento
\citeonline{abntex2modelo}.

\section{Consulte o manual da classe \textsf{abntex2}}

Consulte o manual da classe \textsf{abntex2} \cite{abntex2classe} para uma
referência completa das macros e ambientes disponíveis.

% ---
% Finaliza a parte no bookmark do PDF, para que se inicie o bookmark na raiz
% ---
\bookmarksetup{startatroot}% 
% ---

% ---
% Conclusão
% ---
\section*{Considerações finais}
\addcontentsline{toc}{section}{Considerações finais}

\lipsum[1]

\begin{citacao}
   \lipsum[2]
\end{citacao}

\lipsum[3]

% ----------------------------------------------------------
% ELEMENTOS PÓS-TEXTUAIS
% ----------------------------------------------------------
\postextual

% ---
% Título e resumo em língua estrangeira
% ---

% \twocolumn[    		% INICIO DE ARTIGO EM DUAS COLUNAS

% titulo em inglês
% \titulo{Canonical academic article model with \abnTeX}
% \emptythanks
% \maketitle

% resumo em português
\renewcommand{\resumoname}{Abstract}
\begin{resumoumacoluna}
   \begin{otherlanguage*}{english}
      According to ABNT NBR 6022:2003, an abstract in foreign language is a back
      matter mandatory element.

      \vspace{\onelineskip}

      \noindent
      \textbf{Key-words}: latex. abntex.
   \end{otherlanguage*}
\end{resumoumacoluna}

% ]  				% FIM DE ARTIGO EM DUAS COLUNAS
% ---

% ----------------------------------------------------------
% Referências bibliográficas
% ----------------------------------------------------------
\bibliography{refs}

% ----------------------------------------------------------
% Glossário
% ----------------------------------------------------------
%
% Há diversas soluções prontas para glossário em LaTeX. 
% Consulte o manual do abnTeX2 para obter sugestões.
%
%\glossary

% ----------------------------------------------------------
% Apêndices
% ----------------------------------------------------------

% ---
% Inicia os apêndices
% ---
\begin{apendicesenv}

   % ----------------------------------------------------------
   \chapter{Nullam elementum urna vel imperdiet sodales elit ipsum pharetra ligula
     ac pretium ante justo a nulla curabitur tristique arcu eu metus}
   % ----------------------------------------------------------
   \lipsum[55-57]

\end{apendicesenv}
% ---

% ----------------------------------------------------------
% Anexos
% ----------------------------------------------------------
\cftinserthook{toc}{AAA}
% ---
% Inicia os anexos
% ---
%\anexos
\begin{anexosenv}

   % ---
   \chapter{Cras non urna sed feugiat cum sociis natoque penatibus et magnis dis
     parturient montes nascetur ridiculus mus}
   % ---

   \lipsum[31]

\end{anexosenv}

\end{document}
